% Teacher mark scheme — not for students. Compiler: XeLaTeX.
\documentclass[10pt,a4paper]{article}

\usepackage[margin=16mm,headheight=13pt,headsep=5pt,footskip=9mm]{geometry}
\usepackage{fontspec}
\usepackage[version=4]{mhchem}
\usepackage{booktabs}
\usepackage{array}
\usepackage{enumitem}
\usepackage{fancyhdr}
\usepackage{amsmath,amssymb}

\pagestyle{fancy}
\fancyhf{}
\fancyhead[L]{\small Topic 1 Atomic Structure --- teacher mark scheme}
\fancyhead[R]{\small 34 marks}
\fancyfoot[C]{\thepage}
\renewcommand{\headrulewidth}{0.3pt}

\setlength{\parindent}{0pt}
\setlength{\parskip}{0.12em}
\setlist[enumerate]{leftmargin=1.5em,itemsep=0.10em,topsep=0.12em}

\begin{document}
\begin{center}
{\large\bfseries AS Chemistry --- Topic 1 Atomic Structure (Exam)}\\[0.15em]
{\small Teacher mark scheme · CIE 9701 · 34 marks}
\end{center}

{\small This file is for teachers only. CIE IDs and letters are listed so the paper can be remapped. Student \texttt{main.pdf} has questions only.}

\vspace{0.2em}
{\normalsize\bfseries Section A --- MCQ key}

\vspace{0.1em}
{\small
\begin{tabular}{@{}clll@{}}
\toprule
\textbf{Q} & \textbf{Letter} & \textbf{CIE ID} & \textbf{Paper} \\
\midrule
A1 & \textbf{A} & cie-9701-2021-fm-p12-q1 & 2021 FM P12 Q1 \\
A2 & \textbf{B} & cie-9701-2025-fm-p12-q4 & 2025 FM P12 Q4 \\
A3 & \textbf{A} & cie-9701-2024-mj-p13-q3 & 2024 MJ P13 Q3 \\
A4 & \textbf{C} & cie-9701-2025-mj-p12-q14 & 2025 MJ P12 Q14 \\
A5 & \textbf{D} & cie-9701-2025-on-p11-q40 & 2025 ON P11 Q40 \\
A6 & \textbf{A} & cie-9701-2025-mj-p13-q1 & 2025 MJ P13 Q1 \\
A7 & \textbf{D} & cie-9701-2025-on-p12-q1 & 2025 ON P12 Q1 \\
A8 & \textbf{B} & cie-9701-2021-on-p12-q4 & 2021 ON P12 Q4 \\
A9 & \textbf{C} & cie-9701-2021-mj-p11-q7 & 2021 MJ P11 Q7 \\
A10 & \textbf{C} & cie-9701-2022-mj-p12-q1 & 2022 MJ P12 Q1 \\
A11 & \textbf{C} & cie-9701-2025-mj-p11-q4 & 2025 MJ P11 Q4 \\
A12 & \textbf{D} & cie-9701-2023-on-p11-q2 & 2023 ON P11 Q2 \\
A13 & \textbf{B} & cie-9701-2021-mj-p13-q31 & 2021 MJ P13 Q31 \\
\bottomrule
\end{tabular}}

\vspace{0.35em}
{\normalsize\bfseries B1 · 2025 MJ Paper 21 Question 2 (11 marks)}

\begin{enumerate}[label=\textbf{(\alph*)}, leftmargin=1.8em]
\item [1] atoms (of an element) with same number of protons and different number of neutrons \hfill {\footnotesize cie-9701-2025-mj-p21-q2a}
\item [2] M1 $(5.9\times 54)$, $(91.9\times 56)$, $(2.2\times 57)$. M2 $A_\mathrm{r}=\mathbf{55.9}$ (1 d.p.) \hfill {\footnotesize cie-9701-2025-mj-p21-q2b-i}
\item[] (ii) [2] protons \textbf{26}; nucleons \textbf{56} \hfill {\footnotesize cie-9701-2025-mj-p21-q2b-ii}
\item [1] \textbf{5} \hfill {\footnotesize cie-9701-2025-mj-p21-q2c}
\item [1] \ce{Fe(g) -> Fe+(g) + e-} \hfill {\footnotesize cie-9701-2025-mj-p21-q2d}
\item [4] M1 same electronic configuration / same shielding. M2 same number of protons / same nuclear charge. M3 same nuclear attraction to outer-shell electron. M4 no change / will not be affected. \hfill {\footnotesize cie-9701-2025-mj-p21-q2e}
\end{enumerate}

{\normalsize\bfseries B2 · 2025 MJ Paper 23 Question 1 (10 marks)}

\begin{enumerate}[label=\textbf{(\alph*)}, leftmargin=1.8em]
\item [1] $1s^{2}\,2s^{2}\,2p^{5}$ \hfill {\footnotesize cie-9701-2025-mj-p23-q1a-i}
\item[] (ii) [1] \textbf{2} pairs \hfill {\footnotesize cie-9701-2025-mj-p23-q1a-ii}
\item[] (iii) [1] spherical s orbital (circle) \hfill {\footnotesize cie-9701-2025-mj-p23-q1a-iii}
\item [1] \ce{S(g) -> S+(g) + e-} \hfill {\footnotesize cie-9701-2025-mj-p23-q1b-i}
\item[] (ii) [2] M1 pair of electrons in a 3p orbital (in S) repel. M2 outweighs increase in nuclear charge. \hfill {\footnotesize cie-9701-2025-mj-p23-q1b-ii}
\item[] (iii) [4] M1 \ce{Na+} $<$ \ce{Ne} $<$ \ce{F-}. M2 isoelectronic / same number of electrons. M3 larger species has smaller nuclear charge / fewer protons. M4 smaller nuclear attraction. \hfill {\footnotesize cie-9701-2025-mj-p23-q1b-iii}
\end{enumerate}

\vspace{0.35em}
{\normalsize\bfseries Marks by learning outcome}

\vspace{0.1em}
{\small
\begin{tabular}{@{}lr@{}}
\toprule
\textbf{LO cluster} & \textbf{Marks} \\
\midrule
1.1 particles / nucleon number / electric field / radius & 9 \\
1.2 isotopes & 2 \\
1.2 / 2.1 / 22.2 $A_\mathrm{r}$ from abundances / mass spectrum & 3 \\
1.3 configuration / orbitals / unpaired electrons & 7 \\
1.4 ionisation energy & 13 \\
\midrule
\textbf{Total} & \textbf{34} \\
\bottomrule
\end{tabular}}

\vspace{0.25em}
{\small Quick MCQ string: \textbf{A B A C D · A D B C · C C D B}}

{\small Full exclusion check: \texttt{../TEACHER-KEY.md}.}

\end{document}
